% ============================================================
\chapter{Th\'eor\`eme de Kakutani et Correspondances}
\label{ch:kakutani}
% ============================================================

En 1941, Shizuo Kakutani, jeune math\'ematicien japonais form\'e \`a l'universit\'e d'Osaka, publie un r\'esultat qui semble n'\^etre qu'une extension technique du th\'eor\`eme de Brouwer~: au lieu d'exiger qu'une fonction ait un point fixe, il montre qu'une \emph{correspondance} --- une application qui associe \`a chaque point non pas une valeur unique mais un \emph{ensemble} de valeurs --- poss\`ede elle aussi un point fixe, pourvu que les images soient convexes et que la correspondance soit semi-continue sup\'erieurement. Le r\'esultat aurait pu rester confin\'e aux cercles topologiques. Mais neuf ans plus tard, John Nash l'utilise comme pi\`ece ma\^itresse de sa d\'emonstration de l'existence d'\'equilibres dans les jeux non coop\'eratifs, un travail qui lui vaudra le prix Nobel d'\'economie en 1994. Le th\'eor\`eme de Kakutani est ainsi devenu l'un des piliers de la th\'eorie des jeux et de l'\'economie math\'ematique.

\begin{intuition}
Le th\'eor\`eme de Kakutani g\'en\'eralise le th\'eor\`eme de Brouwer aux applications
multivoques (correspondances). Au lieu d'une fonction $f\colon C \to C$, on consid\`ere
une correspondance $F\colon C \rightrightarrows C$ qui associe \`a chaque point un
\emph{ensemble} de valeurs. Kakutani montre qu'une telle correspondance poss\`ede un point
fixe (un point $x^*$ tel que $x^* \in F(x^*)$) sous des hypoth\`eses de convexit\'e et de
semi-continuit\'e. Ce r\'esultat est fondamental en th\'eorie des jeux, o\`u il permet de
d\'emontrer l'existence d'\'equilibres de Nash.
\end{intuition}

% ------------------------------------------------------------
\section{Correspondances et semi-continuit\'e}
% ------------------------------------------------------------

\begin{definition}[Correspondance (application multivoque)]
\index{correspondance}
Soient $X$ et $Y$ deux espaces topologiques. Une \emph{correspondance} (ou application
multivoque) $F\colon X \rightrightarrows Y$ est une application $F\colon X \to \mathcal{P}(Y)$,
c'est-\`a-dire que pour tout $x \in X$, $F(x)$ est un sous-ensemble de $Y$.
Le \emph{graphe} de $F$ est $\mathrm{Gr}(F) = \{(x,y) \in X \times Y : y \in F(x)\}$.
\end{definition}

\begin{definition}[Semi-continuit\'e sup\'erieure (s.c.s.)]
\index{semi-continuit\'e sup\'erieure}
Une correspondance $F\colon X \rightrightarrows Y$ est \emph{semi-continue sup\'erieurement}
(ou \emph{h\'emi-continue sup\'erieurement}) si pour tout ouvert $V \subseteq Y$,
l'ensemble
\[
  F^+(V) = \{x \in X : F(x) \subseteq V\}
\]
est ouvert dans $X$.
\end{definition}

\begin{definition}[Semi-continuit\'e inf\'erieure (s.c.i.)]
\index{semi-continuit\'e inf\'erieure}
Une correspondance $F\colon X \rightrightarrows Y$ est \emph{semi-continue inf\'erieurement}
si pour tout ouvert $V \subseteq Y$, l'ensemble
\[
  F^-(V) = \{x \in X : F(x) \cap V \neq \emptyset\}
\]
est ouvert dans $X$.
\end{definition}

\begin{proposition}[Caract\'erisation par le graphe ferm\'e]
\label{prop:graphe-ferme}
Soit $F\colon X \rightrightarrows Y$ une correspondance \`a valeurs compactes et $Y$
compact de Hausdorff. Alors $F$ est s.c.s.\ si et seulement si son graphe
$\mathrm{Gr}(F)$ est ferm\'e dans $X \times Y$.
\end{proposition}

\begin{proof}
$(\Rightarrow)$~: Soit $(x_\alpha, y_\alpha)$ un filet dans $\mathrm{Gr}(F)$ convergeant
vers $(x, y)$. Supposons $y \notin F(x)$. Comme $F(x)$ est compact et $Y$ est de
Hausdorff, il existe un ouvert $V$ contenant $F(x)$ tel que $y \notin \overline{V}$.
Par s.c.s., $F^+(V)$ est un voisinage de $x$, donc $F(x_\alpha) \subseteq V$ pour $\alpha$
assez grand, d'o\`u $y_\alpha \in V$. \`A la limite, $y \in \overline{V}$, contradiction.

$(\Leftarrow)$~: Soit $V$ ouvert avec $F(x) \subseteq V$. Si $F^+(V)$ n'\'etait pas un
voisinage de $x$, il existerait un filet $x_\alpha \to x$ avec $F(x_\alpha) \not\subseteq V$.
On trouverait $y_\alpha \in F(x_\alpha) \setminus V$. Par compacit\'e de $Y$, on peut
extraire un sous-filet convergeant vers $y \in Y \setminus V$. Alors $(x,y) \in \mathrm{Gr}(F)$
car le graphe est ferm\'e, d'o\`u $y \in F(x) \subseteq V$, contradiction.
\end{proof}

\begin{remarque}
Pour une fonction univoque $f\colon X \to Y$ (vue comme correspondance $F(x) = \{f(x)\}$),
la s.c.s.\ co\"incide avec la continuit\'e usuelle.
\end{remarque}

% ------------------------------------------------------------
\section{Th\'eor\`eme de Kakutani}
% ------------------------------------------------------------

\begin{theoreme}[Kakutani, 1941]
\label{thm:kakutani}
\index{th\'eor\`eme de Kakutani}
Soit $C$ un sous-ensemble convexe compact non vide de $\R^n$ et
$F\colon C \rightrightarrows C$ une correspondance telle que~:
\begin{enumerate}
  \item pour tout $x \in C$, $F(x)$ est non vide et convexe~;
  \item $F$ est semi-continue sup\'erieurement (graphe ferm\'e).
\end{enumerate}
Alors $F$ poss\`ede un point fixe~: il existe $x^* \in C$ tel que $x^* \in F(x^*)$.
\end{theoreme}

\begin{proof}
Pour chaque $n \in \N^*$, consid\'erons une triangulation simpliciale de $C$ de maille
inf\'erieure \`a $1/n$, dont les sommets sont $v_1^n, \ldots, v_{k_n}^n$. Pour chaque
sommet $v_i^n$, choisissons $w_i^n \in F(v_i^n)$. D\'efinissons l'application affine par
morceaux $f_n\colon C \to C$ par $f_n(v_i^n) = w_i^n$ et extension affine sur chaque
simplexe.

Comme $C$ est convexe compact et $f_n$ est continue, par le th\'eor\`eme de Brouwer, il
existe $x_n \in C$ tel que $f_n(x_n) = x_n$.

Par compacit\'e, $(x_n)$ admet une sous-suite convergente $x_{n_k} \to x^*$. Montrons
que $x^* \in F(x^*)$. Pour tout $\varepsilon > 0$, pour $k$ assez grand, la maille de la
triangulation est inf\'erieure \`a $\varepsilon$, et $x_{n_k}$ appartient \`a un simplexe
dont les sommets $v_i^{n_k}$ v\'erifient $\norm{v_i^{n_k} - x_{n_k}} < \varepsilon$. Comme
$f_{n_k}(x_{n_k}) = x_{n_k}$ est une combinaison convexe des $w_i^{n_k} \in F(v_i^{n_k})$,
et comme $\mathrm{Gr}(F)$ est ferm\'e (Proposition~\ref{prop:graphe-ferme}), on conclut
que $x^* \in F(x^*)$.
\end{proof}

\begin{attention}
L'hypoth\`ese de convexit\'e de $F(x)$ est essentielle. Consid\'erons $C = [0,1]$ et
$F(x) = \{0, 1\} \setminus \{x\}$~: cette correspondance \`a graphe ferm\'e n'a pas de
point fixe car $F(x)$ n'est pas convexe.
\end{attention}

% ------------------------------------------------------------
\section{Extension de Glicksberg--Fan}
% ------------------------------------------------------------

\begin{theoreme}[Glicksberg--Fan]
\label{thm:glicksberg-fan}
Soit $E$ un espace vectoriel topologique localement convexe et $C \subseteq E$ un convexe
compact non vide. Soit $F\colon C \rightrightarrows C$ une correspondance \`a valeurs non
vides, convexes et ferm\'ees, semi-continue sup\'erieurement. Alors $F$ poss\`ede un point
fixe.
\end{theoreme}

\begin{remarque}
Ce th\'eor\`eme g\'en\'eralise Kakutani \`a la dimension infinie, comme le th\'eor\`eme
de Schauder g\'en\'eralise celui de Brouwer.
\end{remarque}

% ------------------------------------------------------------
\section{Application~: \'equilibre de Nash}
% ------------------------------------------------------------

\begin{definition}[Jeu sous forme normale]
\index{jeu!forme normale}
\index{\'equilibre de Nash}
Un \emph{jeu sous forme normale} \`a $n$ joueurs est un triplet
$\Gamma = (N, (S_i)_{i \in N}, (u_i)_{i \in N})$ o\`u~:
\begin{itemize}
  \item $N = \{1, \ldots, n\}$ est l'ensemble des joueurs~;
  \item $S_i$ est l'ensemble des strat\'egies du joueur $i$~;
  \item $u_i\colon S_1 \times \cdots \times S_n \to \R$ est la fonction d'utilit\'e du
    joueur $i$.
\end{itemize}
L'ensemble des strat\'egies mixtes du joueur $i$ est $\Delta(S_i) = \{\sigma_i \in \R^{S_i}
: \sigma_i \geq 0, \sum_{s \in S_i} \sigma_i(s) = 1\}$.
\end{definition}

\begin{definition}[\'Equilibre de Nash]
Un profil de strat\'egies mixtes $\sigma^* = (\sigma_1^*, \ldots, \sigma_n^*)$ est un
\emph{\'equilibre de Nash} si pour tout joueur $i$ et toute strat\'egie mixte
$\sigma_i \in \Delta(S_i)$~:
\[
  u_i(\sigma^*) \geq u_i(\sigma_i, \sigma_{-i}^*)
\]
o\`u $\sigma_{-i}^*$ d\'esigne les strat\'egies des autres joueurs.
\end{definition}

\begin{definition}[Correspondance de meilleure r\'eponse]
La \emph{correspondance de meilleure r\'eponse} du joueur $i$ est
\[
  BR_i(\sigma_{-i}) = \argmax_{\sigma_i \in \Delta(S_i)} u_i(\sigma_i, \sigma_{-i}).
\]
\end{definition}

\begin{theoreme}[Nash, 1950]
\label{thm:nash}
Tout jeu fini sous forme normale (c'est-\`a-dire avec $S_i$ fini pour tout $i$) poss\`ede
au moins un \'equilibre de Nash en strat\'egies mixtes.
\end{theoreme}

\begin{proof}
Posons $C = \prod_{i=1}^n \Delta(S_i)$. C'est un produit de simplexes, donc un convexe
compact non vide de $\R^{\sum \abs{S_i}}$.

D\'efinissons la correspondance de meilleure r\'eponse globale~:
\[
  BR(\sigma) = BR_1(\sigma_{-1}) \times \cdots \times BR_n(\sigma_{-n}).
\]
V\'erifions les hypoth\`eses du th\'eor\`eme de Kakutani~:

\emph{Non-vacuit\'e et convexit\'e}~: $BR_i(\sigma_{-i})$ est non vide (maximum d'une
fonction continue sur un compact) et convexe (car $u_i$ est lin\'eaire en $\sigma_i$~:
si $\sigma_i, \sigma_i'$ maximisent $u_i(\cdot, \sigma_{-i})$, alors toute combinaison
convexe aussi).

\emph{Semi-continuit\'e sup\'erieure}~: le th\'eor\`eme du maximum de Berge assure que
$BR$ est s.c.s., car $u_i$ est continue et $\Delta(S_i)$ ne d\'epend pas de $\sigma_{-i}$.

Par le th\'eor\`eme de Kakutani, $BR$ poss\`ede un point fixe $\sigma^*$, c'est-\`a-dire
$\sigma_i^* \in BR_i(\sigma_{-i}^*)$ pour tout $i$. C'est un \'equilibre de Nash.
\end{proof}

% ------------------------------------------------------------
\section{Th\'eor\`emes minimax}
% ------------------------------------------------------------

\begin{theoreme}[Minimax de von Neumann]
\label{thm:minimax}
\index{th\'eor\`eme minimax}
Soit $A$ une matrice $m \times n$ r\'eelle. Alors~:
\[
  \max_{\sigma \in \Delta_m} \min_{\tau \in \Delta_n} \sigma^T A \tau
  = \min_{\tau \in \Delta_n} \max_{\sigma \in \Delta_m} \sigma^T A \tau.
\]
\end{theoreme}

\begin{proof}
L'in\'egalit\'e $\max\min \leq \min\max$ est toujours vraie. Pour l'\'egalit\'e,
consid\'erons le jeu \`a somme nulle \`a deux joueurs d\'efini par la matrice $A$.
Par le th\'eor\`eme de Nash (ou directement par Kakutani), ce jeu admet un \'equilibre
$(\sigma^*, \tau^*)$. La valeur du jeu est
$v = \sigma^{*T} A \tau^* = \max_\sigma \min_\tau \sigma^T A \tau =
\min_\tau \max_\sigma \sigma^T A \tau$.
\end{proof}

\begin{corollaire}
Tout jeu matriciel \`a somme nulle poss\`ede une valeur, et les deux joueurs ont des
strat\'egies optimales mixtes.
\end{corollaire}

\begin{formules}[title=\textbf{R\'ecapitulatif}]
\begin{center}
\renewcommand{\arraystretch}{1.3}
\begin{tabular}{|l|l|}
\hline
\textbf{Th\'eor\`eme} & \textbf{Hypoth\`eses cl\'es} \\
\hline
Kakutani & $C$ convexe compact $\subseteq \R^n$, $F(x)$ convexe, s.c.s. \\
Glicksberg--Fan & $C$ convexe compact dans evtlc, $F(x)$ convexe ferm\'e, s.c.s. \\
Nash & Jeu fini \`a $n$ joueurs \\
von Neumann & Jeu matriciel \`a somme nulle \\
\hline
\end{tabular}
\end{center}
\end{formules}

% ------------------------------------------------------------
\section{Exercices}
% ------------------------------------------------------------

\begin{exercice}
Montrer que la correspondance $F\colon [0,1] \rightrightarrows [0,1]$ d\'efinie par
$F(x) = [1-x, 1]$ si $x \leq 1/2$ et $F(x) = [0, 1-x]$ si $x > 1/2$ v\'erifie les
hypoth\`eses de Kakutani. Trouver explicitement les points fixes.
\end{exercice}

\begin{exercice}
Soit $C = \{x \in \R^n : \norm{x} \leq 1\}$ la boule unit\'e ferm\'ee. Montrer que la
correspondance $F(x) = \{y \in C : \inner{x}{y} \geq 0\}$ satisfait les hypoth\`eses de
Kakutani et trouver l'ensemble des points fixes.
\end{exercice}

\begin{exercice}[Th\'eor\`eme du maximum de Berge]
\index{th\'eor\`eme de Berge}
Soient $X, Y$ des espaces topologiques, $f\colon X \times Y \to \R$ continue et
$\Gamma\colon X \rightrightarrows Y$ une correspondance s.c.s.\ \`a valeurs compactes non
vides. Montrer que la correspondance de maximiseurs
$M(x) = \argmax_{y \in \Gamma(x)} f(x,y)$ est s.c.s.
\end{exercice}

\begin{exercice}
Trouver les \'equilibres de Nash en strat\'egies mixtes du jeu de Pierre-Feuille-Ciseaux,
d\'efini par la matrice de gains~:
\[
  A = \begin{pmatrix} 0 & -1 & 1 \\ 1 & 0 & -1 \\ -1 & 1 & 0 \end{pmatrix}.
\]
\end{exercice}

\begin{exercice}
Montrer que l'hypoth\`ese de convexit\'e des valeurs $F(x)$ ne peut pas \^etre remplac\'ee
par la connexit\'e. \emph{Indication~: consid\'erer un cercle dans $\R^2$.}
\end{exercice}

\begin{exercice}
Soit $\Gamma$ un jeu \`a deux joueurs o\`u les ensembles de strat\'egies sont des convexes
compacts de $\R^n$ et $\R^m$ respectivement, et les fonctions d'utilit\'e sont continues
et quasi-concaves en la strat\'egie propre de chaque joueur. Montrer l'existence d'un
\'equilibre de Nash en strat\'egies pures.
\end{exercice}
